\documentclass{article}
\usepackage{ctex}
\usepackage{geometry}
\usepackage{xcolor}
\geometry{a4paper}
\title{第四天}
\date{\today}
\author{piyrin}
\begin{document}
\maketitle
\newpage
\begin{abstract}
    滕王阁序练习
\end{abstract}
\newpage
{\color[rgb]{0.4,0.9,0.3}豫章故郡，洪都新府。}星分翼轸，地接衡庐。襟三江而带五湖，控蛮荆而引瓯越。物华天宝，龙光射牛斗之墟；人杰地灵，徐孺下陈蕃之榻。雄州雾列，俊采星驰。台隍枕夷夏之交，宾主尽东南之美。都督阎公之雅望，棨戟遥临；宇文新州之懿范，襜帷暂驻。十旬休假，胜友如云；千里逢迎，高朋满座。腾蛟起凤，孟学士之词宗；紫电青霜，王将军之武库。家君作宰，路出名区；童子何知，躬逢胜饯。

\begin{quotation}
    太史公牛马走司马迁，再拜言。

少卿足下：曩者辱赐书，教以慎于接物，推贤进士为务，意气勤勤恳恳。若望仆不相师，而用流俗人之言，仆非敢如此也。仆虽罢驽，亦尝侧闻长者之遗风矣。顾自以为身残处秽，动而见尤，欲益反损，是以独抑郁而谁与语。谚曰：“谁为为之？孰令听之？”盖钟子期死，伯牙终身不复鼓琴。何则？士为知己者用，女为说己者容。若仆大质已亏缺矣，虽材怀随、和，行若由、夷，终不可以为荣，适足以见笑而自点耳。

书辞宜答，会东从上来，又迫贱事，相见日浅，卒卒无须臾之闲，得竭志意。今少卿抱不测之罪，涉旬月，迫季冬，仆又薄从上雍，恐卒然不可为讳。是仆终已不得舒愤懑以晓左右，则长逝者魂魄私恨无穷。请略陈固陋。阙然久不报，幸勿为过。

仆闻之：修身者，智之符也；爱施者，仁之端也；取予者，义之表也；耻辱者，勇之决也；立名者，行之极也。士有此五者，然后可以托于世，而列于君子之林矣。故祸莫惨于欲利，悲莫痛于伤心，行莫丑于辱先，诟莫大于宫刑。刑余之人，无所比数，非一世也，所从来远矣。昔卫灵公与雍渠同载，孔子适陈；商鞅因景监见，赵良寒心；同子参乘，袁丝变色：自古而耻之！夫以中材之人，事有关于宦竖，莫不伤气，而况于慷慨之士乎！如今朝廷虽乏人，奈何令刀锯之馀荐天下豪俊哉！

且负下未易居，上流多谤议。仆以口语遇此祸，重为乡党所笑，以污辱先人，亦何面目复上父母丘墓乎？虽累百世，垢弥甚耳！是以肠一日而九回，居则忽忽若有所亡，出则不知其所往。每念斯耻，汗未尝不发背沾衣也！身直为闺阁之臣，宁得自引深藏于岩穴邪？故且从俗浮沉，与时俯仰，以通其狂惑。今少卿乃教之以推贤进士，无乃与仆私心刺谬乎？今虽欲自雕琢，曼辞以自解，无益，于俗不信，适足取辱耳。要之死日，然后是非乃定。书不能悉意，故略陈固陋。

仆之先人非有剖符丹书之功，文史星历，近乎卜祝之间，固主上所戏弄，倡优畜之，流俗之所轻也。假令仆伏法受诛，若九牛亡一毛，与蝼蚁何以异？而世又不能与死节者比，特以为智穷罪极，不能自免，卒就死耳。何也？素所自树立使然也。人固有一死，或重于泰山，或轻于鸿毛，用之所趋异也。太上不辱先，其次不辱身，其次不辱理色，其次不辱辞令，其次诬体受辱，其次易服受辱，其次关木索、被箠楚受辱，其次剔毛发、婴金铁受辱，其次毁肌肤、断肢体受辱，最下腐刑极矣！传曰“刑不上大夫”，此言士节不可不勉励也。猛虎在深山，百兽震恐，及在槛阱之中，摇尾而求食，积威约之渐也。故士有画地为牢，势不可入；削木为吏，议不可对，定计于鲜也。今交手足，受木索，暴肌肤，受榜箠，幽于圜墙之中，当此之时，见狱吏则头枪地，视徒隶则心惕息。何者？积威约之势也。及已至是，言不辱者，所谓强颜耳，曷足贵乎！且西伯，伯也，拘于羑里；李斯，相也，具于五刑；淮阴，王也，受械于陈；彭越、张敖，南面称孤，系狱抵罪；绛侯诛诸吕，权倾五伯，囚于请室；魏其，大将也，衣赭衣，关三木；季布为朱家钳奴；灌夫受辱于居室。此人皆身至王侯将相，声闻邻国，及罪至罔加，不能引决自裁，在尘埃之中。古今一体，安在其不辱也？由是言之，勇怯，势也；强弱，形也。审矣，何足怪乎？夫人不能早自裁绳墨之外，以稍陵迟，至于鞭箠之间，乃欲引节，斯不亦远乎！古人所以重施刑于大夫者，殆为此也。夫人情莫不贪生恶死，念父母，顾妻子。至激于义理者不然，乃有所不得已也。今仆不幸，早失父母，无兄弟之亲，独身孤立，少卿视仆于妻子何如哉？且勇者不必死节，怯夫慕义，何处不勉焉！仆虽怯懦，欲苟活，亦颇识去就之分矣，何至自沉溺缧绁之辱哉！且夫臧获婢妾，由能引决，况仆之不得已乎？所以隐忍苟活，幽于粪土之中而不辞者，恨私心有所不尽，鄙陋没世，而文采不表于后世也。

古者富贵而名摩灭，不可胜记，唯倜傥非常之人称焉。盖文王拘而演《周易》；仲尼厄而作《春秋》；屈原放逐，乃赋《离骚》；左丘失明，厥有《国语》；孙子膑脚，《兵法》修列；不韦迁蜀，世传《吕览》；韩非囚秦，《说难》《孤愤》；《诗》三百篇，大底圣贤发愤之所为作也。此人皆意有所郁结，不得通其道，故述往事，思来者。乃如左丘无目，孙子断足，终不可用，退而论书策，以舒其愤，思垂空文以自见。

仆窃不逊，近自托于无能之辞，网罗天下放失旧闻，略考其行事，综其始终，稽其成败兴坏之纪，上计轩辕，下至于兹，为十表，本纪十二，书八章，世家三十，列传七十，凡百三十篇。亦欲以究天人之际，通古今之变，成一家之言。草创未就，会遭此祸，惜其不成，是以就极刑而无愠色。仆诚已著此书，藏之名山，传之其人，通邑大都，则仆偿前辱之责，虽万被戮，岂有悔哉！然此可为智者道，难为俗人言也。

且负下未易居，下流多谤议。仆以口语遇此祸，重为乡党戮笑，污辱先人，亦何面目复上父母之丘墓乎？虽累百世，垢弥甚耳！是以肠一日而九回，居则忽忽若有所亡，出则不知其所往。每念斯耻，汗未尝不发背沾衣也！身直为闺阁之臣，宁得自引深藏于岩穴邪？故且从俗浮沉，与时俯仰，以通其狂惑。今少卿乃教之以推贤进士，无乃与仆私心刺谬乎？今虽欲自雕琢，曼辞以自解，无益，于俗不信，适足取辱耳。要之死日，然后是非乃定。书不能悉意，故略陈固陋。谨再拜。
\end{quotation}

时维九月，序属三秋。潦水尽而寒潭清，烟光凝而暮山紫。俨骖騑于上路，访风景于崇阿；临帝子之长洲，得天人之旧馆。层峦耸翠，上出重霄；飞阁流丹，下临无地。鹤汀凫渚，穷岛屿之萦回；桂殿兰宫，即冈峦之体势。（层峦 一作：层台；即冈 一作：列冈；天人 一作：仙人；飞阁流丹 一作：飞阁翔丹）

披绣闼，俯雕甍，山原旷其盈视，川泽纡其骇瞩。闾阎扑地，钟鸣鼎食之家；舸舰迷津，青雀黄龙之舳。（舳通：轴 迷津 一作：弥津）云销雨霁，彩彻区明。落霞与孤鹜齐飞，秋水共长天一色。渔舟唱晚，响穷彭蠡之滨；雁阵惊寒，声断衡阳之浦。
\newpage    

遥襟甫畅，逸兴遄飞。爽籁发而清风生，纤歌凝而白云遏。睢园绿竹，气凌彭泽之樽；邺水朱华，光照临川之笔。四美具，二难并。穷睇眄于中天，极娱游于暇日。天高地迥，觉宇宙之无穷；兴尽悲来，识盈虚之有数。望长安于日下，目吴会于云间。地势极而南溟深，天柱高而北辰远。关山难越，谁悲失路之人？萍水相逢，尽是他乡之客。怀帝阍而不见，奉宣室以何年？（遥襟甫畅 一作：遥吟俯畅）

嗟乎！时运不齐，命途多舛。冯唐易老，李广难封。所赖君子见机，达人知命。穷且益坚，不坠青云之志。酌贪泉而觉爽，处涸辙以犹欢。北海虽赊，扶摇可接；东隅已逝，桑榆非晚。孟尝高洁，空余报国之情；阮籍猖狂，岂效穷途之哭！（见机 一作：安贫）
\twocolumn
[勃，三尺微命，一介书生。无路请缨，等终军之弱冠；有怀投笔，慕宗悫之长风。舍簪笏于百龄，奉晨昏于万里。非谢家之宝树，接孟氏之芳邻。他日趋庭，叨陪鲤对；今兹捧袂，喜托龙门。胜地不常，盛筵难再；兰亭已矣，梓泽丘墟。临别赠言，幸承恩于伟饯；登高作赋，是所望于群公。敢竭鄙怀，恭疏短引；一言均赋，四韵俱成。请洒潘江，各倾陆海云尔：

滕王高阁临江渚，佩玉鸣鸾罢歌舞。 画栋朝飞南浦云，珠帘暮卷西山雨。 闲云潭影日悠悠，物换星移几度秋。 槛外长江空自流。

这篇骈文以生动的笔触描绘了滕王阁的壮丽景色和宴会的盛况，同时抒发了作者的个人抱负和怀才不遇的愤懑之情。全文对仗工整，音韵和谐，辞藻华丽，是骈文中的佳作。
这里是汉代的豫章郡城，如今是新设立的洪州府，天上的方位属于翼、轸两星宿的分野，地上的位置连结着衡山和庐山。以三江为衣襟，以五湖为衣带，控制着楚地，连接着瓯越。这里物产的华美，有如天降之宝，其光彩上冲牛斗之宿。这里的土地有灵秀之气，陈蕃专为徐孺设下几榻。洪州境内的建筑如云雾排列，有才能的人士如流星一般奔驰驱走。城池坐落在中原与南夷的交界之处，宾客与主人包括了东南地区最优秀的人物。都督阎公，享有崇高的名望，远道来到洪州坐镇，宇文州牧，是美德的楷模，赴任途中在此暂留。每逢十日一旬的假期，来了很多的良友，迎接远客，高贵的朋友坐满了席位。文词宗主孟学士所作文章就像腾起的蛟龙、飞舞的彩凤；王将军的兵器库中，藏有像紫电、青霜这样锋利的宝剑。由于父亲在交趾做县令，我在探亲途中经过这个著名的地方。我年幼无知，竟有幸亲身参加了这次盛大的宴会。历史

　　正当深秋九月之时，雨后的积水消尽，寒凉的潭水清澈，天空凝结着淡淡的云烟，暮霭中山峦呈现一片紫色。在高高的山路上驾着马车，在崇山峻岭中访求风景。来到昔日帝子的长洲，找到仙人居住过的宫殿。这里山峦重叠，青翠的山峰耸入云霄。凌空的楼阁，红色的阁道犹如飞翔在天，从阁上看不到地面。仙鹤野鸭栖止的水边平地和水中小洲，极尽岛屿的纡曲回环之势；华丽威严的宫殿，依凭起伏的山峦而建。

　　推开雕花精美的阁门，俯视彩饰的屋脊，山峰平原尽收眼底，河流迂回的令人惊讶。遍地是里巷宅舍，许多钟鸣鼎食的富贵人家。舸舰塞满了渡口，尽是雕上了青雀黄龙花纹的大船。云消雨停，阳光普照，天空晴朗；落日映射下的彩霞与孤单的野鸭一齐飞翔，秋天的江水和辽阔的天空连成一片，浑然一色。傍晚时分，渔夫在渔船上歌唱，那歌声响彻彭蠡湖滨；深秋时节，雁群感到寒意而发出惊叫，哀鸣声一直持续到衡阳的水滨。诗歌]

　　放眼远望，胸襟顿时感到舒畅，超逸的兴致立即兴起。排箫的音响引来徐徐清风，柔缓的歌声吸引住飘动的白云。今日盛宴好比当年梁园雅集，大家酒量也胜过陶渊明。参加宴会的文人学士，就像当年的曹植，写出“朱华冒绿池”一般的美丽 诗句，其风流文采映照着谢灵运的诗笔。音乐与饮食、文章和言语这四种美好的事物都已经齐备，贤主、嘉宾这两个难得的条件也凑合在一起了。向天空中极目远眺，在假日里尽情欢娱。苍天高远，大地寥廓，令人感到宇宙的无穷无尽。欢乐逝去，悲哀袭来，意识到万事万物的的消长兴衰是有定数的。远望长安沉落到夕阳之下，遥看吴郡隐现在云雾之间。地理形势极为偏远，南方大海特别幽深，昆仑山上天柱高耸，缈缈夜空北极远悬。关山重重难以越过，有谁同情我这不得志的人？偶然相逢，满座都是他乡的客人。怀念着君王的宫门，但却不被召见，什么时候才能像贾谊那样到宣室侍奉君王呢？

gfd 
\begin{minipage}[t]{8em}
呵，各人的时机不同，\\
人生的命运多有不顺。\\
冯唐容易衰老，李广立功无数却难得封侯\\
使贾谊这样有才华的人屈居于长沙，并不是当时没有圣明的君主，使梁鸿逃匿到齐鲁海滨，不是在政治昌明的时代吗？只不过由于君子能了解时机，通达的人知道自己的命运罢了。
\end{minipage}
gfd 


fdf\rule{10em}{1pt}
　　年岁虽老而心犹壮，怎能在白头时改变心情？遭遇穷困而意志更加坚定，在任何情况下也不放弃自己的凌云之志。即使喝了贪泉的水，也觉得清爽可口，并不滋生贪心；即使像鲋鱼处于即将干涸的车辙中，依然开朗愉快。北海虽然遥远，乘着大风仍然可以到达；晨光虽已逝去，珍惜黄昏却为时不晚。孟尝心性高洁，但白白地怀抱着报国的热情，阮籍为人放纵不羁，我怎能学他走到穷途就哭泣的行为呢！

　　我地位卑微，只是一介书生。虽然和终军年龄相等，却没有报国的机会。像班超那样有投笔从戎的豪情，也有宗悫“乘风破浪”的壮志。如今我抛弃了一生的功名，不远万里去朝夕侍奉父亲。虽然不是谢玄那样的人才，但也和许多贤德之士相交往。过些日子，我将到父亲身边，一定要像孔鲤那样接受父亲的教诲；而今天我能谒见阎公受到接待，高兴得如同登上龙门一样。假如碰不上杨得意那样引荐的人，就只有抚拍着自己的文章而自我叹惜。既然已经遇到了钟子期，就弹奏一曲《流水》又有什么羞愧呢？历史


\setlength{\fboxrule}{5pt}
\setlength{\fboxsep}{2em}
\fbox{呵！名胜之地不能常存，盛大的宴会难以再逢。兰亭集会的盛况已成陈迹，石崇的梓泽也变成了废墟。承蒙这个宴会的恩赐，让我临别时作了这一篇序文，至于登高作赋，这只有指望在座诸公了。我只是冒昧地尽我微薄的心意，作了短短的引言。我的一首四韵小 诗也已写成。请各位像潘岳、陆机那样，展现江海般的文才吧：

\begin{verse}
    巍峨高耸的滕王阁俯临着江心的沙洲，想当初佩玉、鸾铃鸣响的豪华歌舞已经停止了。
　　早晨，南浦轻云掠过滕王阁的画栋；傍晚时分，西山烟雨卷起滕王阁的珠帘。
　　悠闲的彩云影子倒映在江水

\end{verse}


}

\rlap{这是}方式反对法中，整天悠悠然地漂浮着；时光易逝，人事变迁，不知已经度过几个春秋。
　　昔日游赏于高阁中的滕王如今已不知哪里去了，只有那栏杆外的滔滔江水空自向远方奔流。诗歌
\marginpar{\footnotesize fjskfds}
\footnote {滕王阁序}
\begin{verbatim}
    #include<iostream>
    using namespace std;
    int main()
    {
        cout<<"hello world";
        return 0;
    }

\end{verbatim}
\end{document}